\chapter{Conclusiones y Trabajo Futuro}
\label{chap:conclusiones}

\section{Conclusiones Generales}
\label{sec:conclusiones_generales}

Este Trabajo de Fin de Máster ha abordado el problema de la selección de carteras financieras bajo restricciones
discretas de cardinalidad mediante un algoritmo cuántico híbrido basado en QAOA, comparando su comportamiento
frente a un solucionador exacto clásico (Gurobi) y frente a una metaheurística de referencia industrial (Recocido
Simulado). El propósito del proyecto era responder a una pregunta concreta: \emph{¿puede un ansatz variacional
cuántico, diseñado para respetar restricciones combinatorias duras, resolver el problema de factibilidad del QAOA
estándar sin sacrificar una calidad de solución razonable frente a los métodos clásicos establecidos?}

Los resultados del Capítulo~\ref{ch:resultados} responden a esta pregunta de forma matizada. La sustitución de
las penalizaciones QUBO blandas por el confinamiento geométrico del \emph{XY-Mixer} sobre el estado de Dicke
resuelve de forma estructural el problema de viabilidad: mientras que el QAOA estándar degrada su tasa de
factibilidad hasta un $51\text{--}52\,\%$ en $N=18\text{--}20$, la arquitectura XY-QAOA mantiene una viabilidad
matemática del $100\,\%$ en todas las instancias evaluadas, por construcción y no por ajuste estadístico
(Sección~\ref{sec:factibilidad}). Es, con diferencia, la evidencia más sólida de todo el estudio.

En calidad de solución (Optimization GAP), el panorama cambia: bajo el régimen de un único reinicio evaluado en
este trabajo, XY-QAOA Regularizado mejora de forma clara y consistente a Standard QAOA en los ocho valores de
$N$ estudiados, pero \textbf{no supera a Simulated Annealing}, que se mantiene en una meseta de
$4\text{--}4{,}5\,\%$ frente al $11\text{--}22\,\%$ de XY-QAOA Regularizado (Sección~\ref{sec:gap-vs-n}). Esta
diferencia es estadísticamente significativa en los ocho valores de $N$ (test de Wilcoxon pareado por semilla,
$p<0{,}01$ en todos los casos), por lo que no se debe al azar de muestreo entre las 12 semillas evaluadas. En
robustez financiera fuera de muestra, el Sharpe OOS de XY-QAOA Regularizado se sitúa entre el $76\,\%$ y el
$95\,\%$ del valor de Gurobi OOS según $N$, sin diferencia estadísticamente significativa en ningún punto
($p>0{,}05$ en los ocho valores de $N$); ambos enfoques son, por tanto, financieramente \textbf{comparables}
fuera de muestra.

En conjunto, estas evidencias respaldan una tesis más acotada que la planteada inicialmente, pero más defendible:
en el régimen NISQ actual, bajo simulación clásica exacta y con un presupuesto mínimo de reinicios, un algoritmo
cuántico variacional bien diseñado \textbf{resuelve de forma determinista el problema de factibilidad} que
aqueja al QAOA estándar, y alcanza una robustez financiera fuera de muestra \textbf{comparable} a la de un
solucionador clásico exacto --pero no supera en calidad de solución (GAP) a una metaheurística clásica simple
como Simulated Annealing. El valor diferencial de este trabajo es, por tanto, la garantía estructural de
factibilidad, no una ventaja de optimalidad ni de generalización financiera.

\section{Grado de Cumplimiento de los Objetivos e Hipótesis}
\label{sec:cumplimiento_objetivos}

Contrastando los resultados del Capítulo~\ref{ch:resultados} con los objetivos específicos de la
Sección~\ref{sec:objetivos_especificos}: OE1 (equivalencia QUBO--Ising) queda validado por la verificación de la
Sección~\ref{sec:verificacion}, con coincidencia exacta en el $100\,\%$ de las 15 combinaciones evaluadas. OE2
(factibilidad estructural) se cumple con holgura: $100\,\%$ de viabilidad en todo el rango $N=6$ a $N=20$, frente
a la degradación progresiva de Standard QAOA (Sección~\ref{sec:factibilidad}).

OE3 (efecto estabilizador de TQA y Ridge $L_2$) se cumple solo \textbf{parcialmente}. Aislando el efecto puro de
TQA --comparando XY-QAOA con inicialización aleatoria frente a XY-QAOA Regularizado, ambos con el mismo
mezclador $XY$-- el test de Wilcoxon pareado no encuentra diferencia significativa en el GAP medio en ninguno de
los ocho valores de $N$ ($p$ entre $0{,}07$ y $0{,}90$), y la reducción de varianza atribuible a TQA es
inconsistente entre valores de $N$ (entre $-11{,}2\,\%$ y $+34{,}1\,\%$, media $6{,}4\,\%$). La regularización
Ridge $L_2$ no aporta mejora aislada alguna (Sección~\ref{sec:limitaciones}). El objetivo se cumple en el
sentido de que se ha evaluado rigurosamente el efecto de ambos mecanismos, pero la evidencia no respalda que
ninguno de los dos estabilice, por sí solo, el bucle variacional de forma estadísticamente significativa en el
rango estudiado.

OE4 (contraste frente a solucionadores clásicos) se cumple en el sentido de que se ha hecho la comparación
sistemática solicitada, con un resultado claro aunque no el anticipado: XY-QAOA Regularizado no iguala a
Simulated Annealing en GAP (diferencia significativa, $p<0{,}01$ en los ocho valores de $N$), pero sí resulta
financieramente comparable a Gurobi en Sharpe OOS (diferencia no significativa, $p>0{,}05$ en los ocho valores
de $N$), y el umbral a partir del cual la simulación clásica deja de ser práctica se sitúa entre $N=14$ y $N=16$
(Sección~\ref{sec:tiempo-vs-n}).

En conjunto, se considera cumplido el objetivo general en su vertiente de diseño, implementación y evaluación
sistemática del algoritmo propuesto. El grado de competitividad frente a los métodos clásicos debe matizarse:
completo en factibilidad, comparable en robustez financiera fuera de muestra, y no alcanzado en calidad de
solución frente a una metaheurística clásica simple bajo este diseño experimental concreto.

\subsection{Contraste de las Hipótesis de Trabajo}

La Tabla~\ref{tab:hipotesis} resume el grado de confirmación de las hipótesis de la Sección~\ref{sec:hipotesis},
a la luz de la evidencia estadística del Capítulo~\ref{ch:resultados}.

\begin{table}[htbp]
\centering
\small
\caption{Grado de confirmación de las hipótesis de trabajo a la luz de la evidencia experimental.}
\label{tab:hipotesis}
\begin{tabular}{p{2.6cm}p{7.2cm}p{2.8cm}}
\toprule
\textbf{Hipótesis} & \textbf{Evidencia} & \textbf{Resultado} \\
\midrule
H1: el XY-Mixer mejora la factibilidad &
$100\,\%$ de factibilidad en XY-QAOA/XY-QAOA Reg. en todo el rango, frente a una caída de
Standard QAOA hasta $\sim\!52\,\%$ en $N=18$--$20$ (Sección~\ref{sec:factibilidad}) &
\textbf{Confirmada} \\
\addlinespace
H2: TQA estabiliza la optimización &
Aislando el efecto de TQA frente a inicialización aleatoria (mismo mezclador $XY$), sin
diferencia significativa en el GAP medio ($p>0{,}07$ en los 8 valores de $N$) ni reducción
consistente de varianza &
\textbf{No confirmada} \\
\addlinespace
H3: Ridge $L_2$ aporta mejora adicional &
La ablación de $\alpha$ no muestra mejora consistente; el término queda desactivado en la
configuración final (Sección~\ref{sec:limitaciones}) &
\textbf{No confirmada} \\
\addlinespace
H4: competitividad frente a clásicos &
GAP significativamente peor que Simulated Annealing ($p<0{,}01$ en los 8 valores de $N$);
Sharpe OOS sin diferencia significativa frente a Gurobi ($p>0{,}05$) &
\textbf{Parcialmente confirmada} (competitivo en Sharpe, no en GAP) \\
\addlinespace
H5: límite de expresividad a $p$ fijo &
Repunte del GAP en $N=18$--$20$; el escalado en $p$ (Sección~\ref{sec:efecto-p}) sugiere
margen de mejora al aumentar la profundidad &
\textbf{Confirmada} \\
\bottomrule
\end{tabular}
\end{table}

\section{Limitaciones del Trabajo}
\label{sec:limitaciones}

A pesar de los resultados obtenidos, conviene reconocer las limitaciones del estudio:

\begin{itemize}
    \item \textbf{Simulación sin ruido físico.} Todos los experimentos se ejecutaron mediante simulación exacta
    de vector de estado (\emph{statevector}), sin modelar canales de ruido de hardware real. Los resultados de
    viabilidad y GAP representan, por tanto, una cota superior de rendimiento frente a un procesador cuántico
    físico actual.
    \item \textbf{XY-QAOA Regularizado no supera a Simulated Annealing en calidad de solución.} Bajo el diseño
    experimental de este trabajo (un único reinicio, simulación exacta), el GAP medio del enfoque cuántico es
    significativamente peor que el de una metaheurística clásica simple en todo el rango de $N$. Esta limitación
    es central para interpretar el alcance de la contribución: el valor demostrado reside en la garantía
    estructural de factibilidad, no en una ventaja de optimalidad.
    \item \textbf{Contribución de TQA y de la regularización Ridge $L_2$ no confirmada estadísticamente.} Como se
    detalla en la Sección~\ref{sec:cumplimiento_objetivos}, ni la inicialización TQA ni la penalización Ridge
    $L_2$ muestran, por sí solas, una mejora estadísticamente significativa sobre el GAP medio o su variabilidad
    en el rango de $N$ estudiado.
    \item \textbf{Límite de expresividad a profundidad fija.} El repunte del GAP en $N=18$--$20$ confirma que,
    para una profundidad de circuito $p$ fija, el ansatz agota su capacidad de representación al crecer $N$.
    \item \textbf{Coste computacional de la simulación clásica.} El tiempo de ejecución escala mucho peor que los
    solucionadores clásicos de referencia, limitando la escala de universos de activos evaluables y alejando,
    por ahora, la viabilidad de un despliegue en tiempo real.
    \item \textbf{Alcance del modelo financiero.} El estudio se restringe a un modelo estático de media-varianza
    de un único periodo sobre renta variable (S\&P 500 / IBEX), sin costes de transacción ni rebalanceo dinámico.
    \item \textbf{Tamaño muestral moderado.} Cada punto se apoya en $n=12$ semillas, lo que produce intervalos de
    confianza relativamente amplios, en particular para Standard QAOA en $N$ elevado (Sección~\ref{sec:limitaciones-estadisticas}).
\end{itemize}

\section{Líneas Futuras de Trabajo}
\label{sec:trabajo_futuro}

A partir de las limitaciones anteriores, se identifican las siguientes líneas de continuación, de menor a mayor
alcance:

\begin{enumerate}
    \item \textbf{Validación en hardware cuántico real.} Ejecutar el circuito XY-QAOA en plataformas NISQ
    físicas para contrastar la degradación de viabilidad y GAP frente a la simulación exacta.
    \item \textbf{Ampliación del número de reinicios y de semillas.} Dado que Standard QAOA y, en menor medida,
    XY-QAOA Regularizado muestran varianza inter-semilla elevada con un único reinicio, evaluar si un
    presupuesto de reinicios mayor (más realista en emuladores clásicos que en hardware NISQ) reduce la brecha
    de GAP frente a Simulated Annealing.
    \item \textbf{Búsqueda de hiperparámetros más amplia para TQA y Ridge.} Como no se ha confirmado
    estadísticamente una mejora de estos mecanismos en el rango estudiado, explorar con optimización bayesiana
    la rampa temporal de TQA y el coeficiente $\alpha$ de Ridge, para ver si existe una región del espacio de
    hiperparámetros donde sí aporten una mejora consistente.
    \item \textbf{Estrategias de profundidad adaptativa.} Investigar esquemas de \emph{warm-start} o de
    crecimiento incremental de $p$ (p.\,ej. ADAPT-QAOA) para compensar el límite de expresividad detectado.
    \item \textbf{Extensión del modelo financiero.} Incorporar rebalanceo dinámico multiperiodo, costes de
    transacción y restricciones sectoriales o de liquidez.
    \item \textbf{Ampliación del universo de activos y regímenes de mercado.} Evaluar el algoritmo sobre otras
    clases de activos y bajo distintos regímenes de mercado.
    \item \textbf{Exploración de arquitecturas alternativas.} Comparar el XY-Mixer con otros esquemas de
    preservación de restricciones (p.\,ej. \emph{Grover-mixer}, QAOA+ multiángulo) que puedan cerrar la brecha de
    GAP frente a Simulated Annealing.
    \item \textbf{Análisis de coste energético a escala.} Extender el análisis de huella de carbono a escenarios
    de despliegue en QPU real.
\end{enumerate}

En definitiva, este trabajo demuestra que la preservación estructural de restricciones mediante circuitos
cuánticos especializados resuelve de forma robusta y verificable el problema de factibilidad del QAOA estándar,
y ofrece una robustez financiera fuera de muestra comparable a la de un solucionador clásico exacto. Cerrar la
brecha de calidad de solución frente a metaheurísticas clásicas simples queda como el reto más relevante para
que el enfoque sea competitivo en optimalidad, y no solo en factibilidad, en próximas iteraciones de este
trabajo.